% !PRINTED 510
\weberpage{510}

3. \gesperrt{Ist $\alpha$ theilbar durch $\beta$, und $\beta$ theilbar durch $\gamma$,
so ist auch $\alpha$ theilbar durch $\gamma$.}

Denn nach der Voraussetzung giebt es zwei ganze Functionale $\varkappa$, $\lambda$, die
den Bedingungen $\alpha = \varkappa\beta$, $\beta = \lambda\gamma$ genügen. Demnach ist
auch $\alpha = \varkappa\lambda\gamma$, und da $\varkappa\lambda$ ganz ist, so ist
$\alpha$ durch $\gamma$ theilbar.

Ein Product $\beta\gamma$ zweier ganzer Functionale ist sowohl durch $\beta$ als durch
$\gamma$ theilbar, und folglich ist, wenn $\beta$ durch $\alpha$ theilbar ist, auch
$\beta\gamma$ durch $\alpha$ theilbar. Aus diesen Sätzen ergiebt sich:

4. \gesperrt{Sind $\alpha$, $\beta$, \ldots\ durch $\delta$ theilbar, und sind $\xi$,
$\eta$, \ldots\ beliebige ganze Functionale, so ist $\xi\alpha + \eta\beta + \ldots$
durch $\delta$ theilbar.}

Selbstverständlich erstrecken sich diese Definitionen und Sätze auch auf den Fall, dass
Zahlen an die Stelle von Functionalen treten, und so erhalten wir die Theilbarkeit
ganzer Zahlen.

5. \gesperrt{Zwei ganze Functionale $\alpha$, $\beta$, die gegenseitig durch einander
theilbar sind, heissen associirt.}

6. \gesperrt{Ein mit der natürlichen Zahl 1 associirtes ganzes Functional $\varepsilon$,
d.\,h. jeder Theiler der Zahl 1, heisst eine Einheit.}

Je nachdem $\varepsilon$ ein Functional oder eine Zahl ist, ist $\varepsilon$ eine
\gesperrt{functionale} oder eine \gesperrt{numerische} Einheit.

Im Körper $R$ der rationalen Zahlen sind als functionale Einheiten die primitiven
ganzen Functionen und die Quotienten von zweien unter ihnen anzusehen. Numerische
Einheiten giebt es in $R$ nur zwei, nämlich $+\,1$ und $-\,1$.

Ueber die hierdurch eingeführten Begriffe, die in enger gegenseitiger Beziehung
stehen, leiten wir eine Reihe von Sätzen ab.

7. \gesperrt{Sind $\alpha$, $\beta$ associirte Functionale, so sind die Quotienten
$\beta : \alpha$ und $\alpha : \beta$ Einheiten.}

Denn setzen wir $\beta = \alpha\varepsilon$, so ist $\varepsilon$ ein ganzes
Functional und $1 : \varepsilon = \alpha : \beta$ gleichfalls ganz; also ist
$\varepsilon$ ein Theiler der Zahl 1, d.\,h. $\varepsilon$ ist eine Einheit.

8. \gesperrt{Ist $\alpha$ ein ganzes Functional und $\varepsilon$ eine Einheit, so
sind $\alpha$ und $\alpha\varepsilon$ associirt.}

Dies folgt unmittelbar aus der Definition; denn $\alpha : \alpha\varepsilon =
1 : \varepsilon$ und $\alpha\varepsilon : \alpha = \varepsilon$ sind beides ganze
Functionale.
